\documentclass[11pt]{article}
\usepackage[a4paper,margin=1.05in]{geometry}
\usepackage[T1]{fontenc}
\usepackage[utf8]{inputenc}
\usepackage{lmodern}
\usepackage{amsmath,amssymb,amsthm,mathtools}
\usepackage{enumitem}
\usepackage[hidelinks]{hyperref}
\usepackage{microtype}

\newtheorem{theorem}{Theorem}[section]
\newtheorem{proposition}[theorem]{Proposition}
\newtheorem{lemma}[theorem]{Lemma}
\newtheorem{corollary}[theorem]{Corollary}
\newtheorem{definition}[theorem]{Definition}
\newtheorem{remark}[theorem]{Remark}
\newtheorem{example}[theorem]{Example}

\DeclareMathOperator{\im}{im}
\DeclareMathOperator{\rank}{rank}
\DeclareMathOperator{\spanC}{span}
\DeclareMathOperator{\charac}{char}
\DeclareMathOperator{\Supp}{Supp}

\title{Filtered Equivariant Membership and Static Algebraic Proof Degree}
\author{Luca Blanchi}
\date{June 2026}

\begin{document}
\maketitle

\begin{abstract}
We develop a presentation-sensitive, filtered, and equivariant framework for static algebraic proof degree. Let \(k\) be a field, let \(R=k[x_1,\ldots,x_n]\), and let
\[
\mathbf F=(f_1,\ldots,f_m)
\]
be a finite polynomial presentation. For a consequence \(h\in(f_1,\ldots,f_m)\), its static derivation degree is
\[
\delta_{\mathbf F}(h)
=
\min\left\{D:
h=\sum_i g_i f_i,\ \deg(g_i f_i)\le D\ \text{for all }i
\right\}.
\]
If \(J_{\mathbf F}=(f_1^h,\ldots,f_m^h)\subset k[x_0,x_1,\ldots,x_n]\) is the ideal generated by the homogenizations of the chosen generators, then
\[
\delta_{\mathbf F}(h)
=
\min\{D\ge \deg h:x_0^{D-\deg h}h^h\in J_{\mathbf F}\}.
\]
In particular, the static Hilbert--Nullstellensatz refutation degree is
\[
\delta_{\mathbf F}(1)=\min\{D:x_0^D\in J_{\mathbf F}\}.
\]

The point is not the homogenization identity alone, which is close to standard homogenized Polynomial Calculus principles. The point is the filtered interpretation it exposes. The ideal \(J_{\mathbf F}\) is the homogenized ideal of the chosen presentation, not the saturated homogenization of the affine ideal, and bounded static derivability is controlled by the \(x_0\)-torsion of
\[
(J_{\mathbf F}:x_0^\infty)/J_{\mathbf F}.
\]
Over finite domains, static refutation degree becomes exact filtered membership in the coordinate algebra of the domain. For a single equation \(f=0\) with no zeros on a finite degree-compatible domain \(X\),
\[
\delta=\deg f+\deg_X(f^{-1}).
\]
On the Boolean cube this gives exact formulas in terms of the algebraic degree of reciprocal Boolean functions; generically a degree-\(d\) Boolean contradiction has static degree \(n+d\), and the cardinality contradiction
\[
x_i^2-x_i=0,\qquad x_1+\cdots+x_n-a=0,\qquad a\notin\{0,\ldots,n\}
\]
has exact degree \(n+1\).

Finally, if a linearly reductive group acts on the presentation and fixes the homogenizing coordinate, invariant targets are tested exactly in invariant Macaulay maps, while general targets are controlled by their isotypic components. Symmetric Boolean systems reduce exactly to univariate filtered membership in
\[
k[t]/\prod_{a=0}^n(t-a),
\]
and block-symmetric systems reduce to filtered membership on finite grids of orbit parameters.
\end{abstract}

\section{Static Algebraic Degree}

Let \(k\) be a field and let
\[
R=k[x_1,\ldots,x_n].
\]
A finite ordered list
\[
\mathbf F=(f_1,\ldots,f_m)
\]
of nonzero polynomials will be called a presentation. It generates the ideal
\[
I_{\mathbf F}=(f_1,\ldots,f_m)\subset R.
\]
The chosen generating set matters: two presentations of the same ideal can have different bounded-degree derivations.

\begin{definition}[Static derivation degree]
For \(h\in R\), define
\[
\delta_{\mathbf F}(h)
=
\min
\left\{
D:
h=\sum_{i=1}^m g_i f_i,\quad
g_i\in R,\quad
\deg(g_i f_i)\le D\ \text{for all }i
\right\}.
\]
If no such representation exists, set \(\delta_{\mathbf F}(h)=\infty\). For \(h=0\), set \(\delta_{\mathbf F}(0)=0\).
\end{definition}

When \(h=1\), this is the minimum static Hilbert--Nullstellensatz refutation degree of the presentation, provided \(1\in I_{\mathbf F}\). This is a degree invariant of static ideal-membership identities. It is not a size measure for dynamic Polynomial Calculus proofs or for algebraic circuits.

\section{Homogenized Presentation Ideals}

Let
\[
S=k[x_0,x_1,\ldots,x_n].
\]
For \(0\ne p\in R\), let \(p^h\in S\) denote its homogenization to total degree \(\deg p\):
\[
p^h=x_0^{\deg p}p(x_1/x_0,\ldots,x_n/x_0).
\]
Given a presentation \(\mathbf F=(f_1,\ldots,f_m)\), define
\[
J_{\mathbf F}=(f_1^h,\ldots,f_m^h)\subset S.
\]
This is the ideal generated by the homogenizations of the chosen generators. It is generally smaller than the saturated homogenization of \(I_{\mathbf F}\), and this distinction is exactly what records the degree cost of the chosen presentation.

For a homogeneous ideal \(J\subset S\) and a homogeneous polynomial \(H\in S_r\), define
\[
\nu_{x_0}(J;H)
=
\min\{D\ge r:x_0^{D-r}H\in J\},
\]
with value \(\infty\) if no such \(D\) exists. For \(H=1\), write
\[
\nu_{x_0}(J)=\min\{D:x_0^D\in J\}.
\]

\begin{theorem}[Homogenized membership formula]
For every nonzero \(h\in R\),
\[
\delta_{\mathbf F}(h)
=
\nu_{x_0}(J_{\mathbf F};h^h)
=
\min\{D\ge \deg h:x_0^{D-\deg h}h^h\in J_{\mathbf F}\}.
\]
In particular, if \(1\in I_{\mathbf F}\), then
\[
\delta_{\mathbf F}(1)=\nu_{x_0}(J_{\mathbf F})
=
\min\{D:x_0^D\in J_{\mathbf F}\}.
\]
\end{theorem}

\begin{proof}
Let \(r=\deg h\). Suppose first that
\[
h=\sum_i g_i f_i
\]
and \(\deg(g_i f_i)\le D\) for all \(i\). Homogenizing the identity to total degree \(D\) gives
\[
x_0^{D-r}h^h
=
\sum_i x_0^{D-\deg(g_i f_i)}g_i^h f_i^h,
\]
so \(x_0^{D-r}h^h\in J_{\mathbf F}\).

Conversely, suppose \(x_0^{D-r}h^h\in J_{\mathbf F}\). Since \(J_{\mathbf F}\) is homogeneous, its degree-\(D\) part gives homogeneous polynomials \(G_i\), with \(G_i=0\) if \(D<\deg f_i\), such that
\[
x_0^{D-r}h^h=\sum_i G_i f_i^h,\qquad \deg G_i=D-\deg f_i
\]
whenever \(G_i\ne0\). Setting \(x_0=1\), we obtain
\[
h=\sum_i G_i(1,x_1,\ldots,x_n)f_i
\]
and each product has degree at most \(D\). Taking minima proves the formula. The case \(h=1\) is the specialization \(1^h=1\).
\end{proof}

\section{\texorpdfstring{\(x_0\)}{x0}-Torsion and Saturation Defect}

Let \(J=J_{\mathbf F}\). Define
\[
J:x_0^\infty=\{H\in S:\exists e\ge0,\ x_0^eH\in J\}
\]
and the saturation-defect module
\[
T_{\mathbf F}=(J_{\mathbf F}:x_0^\infty)/J_{\mathbf F}.
\]
If \(h\in I_{\mathbf F}\), then \(h^h\in J_{\mathbf F}:x_0^\infty\), because any affine expression for \(h\) homogenizes after multiplying by a sufficiently high power of \(x_0\).

\begin{proposition}[Torsion order]
For every nonzero \(h\in I_{\mathbf F}\),
\[
\delta_{\mathbf F}(h)-\deg h
=
\min\{e\ge0:x_0^e h^h\in J_{\mathbf F}\}.
\]
Equivalently, \(\delta_{\mathbf F}(h)-\deg h\) is the \(x_0\)-torsion order of the class of \(h^h\) in \(T_{\mathbf F}\).
\end{proposition}

\begin{proof}
Put \(e=D-\deg h\) in the homogenized membership formula.
\end{proof}

Thus static degree measures how long a homogeneous consequence survives in the \(x_0\)-torsion defect of the chosen homogenized presentation.

\section{Degeneration Bounds}

The preceding formulation gives simple monotonicity statements under homogeneous flat specialization.

\begin{proposition}[Flat specialization]
Let \(A\) be an integral \(k\)-algebra and let
\[
\mathcal J\subset S\otimes_k A
\]
be a homogeneous ideal such that \((S\otimes_k A)/\mathcal J\) is flat over \(A\). Let \(J_\eta\) be the generic fiber and \(J_s\) a specialization. Then
\[
\nu_{x_0}(J_\eta)\ge \nu_{x_0}(J_s).
\]
\end{proposition}

\begin{proof}
For fixed \(D\), the degree-\(D\) piece of \((S\otimes_k A)/\mathcal J\) is flat, hence torsion-free over the integral domain \(A\). If the class of \(x_0^D\) vanishes generically, torsion-freeness forces it to vanish before specialization, hence in every fiber. Taking the least such \(D\) gives the claim.
\end{proof}

\begin{corollary}[Initial ideals]
For every homogeneous ideal \(J\subset S\) and every homogeneous initial ideal \(\operatorname{in}(J)\),
\[
\nu_{x_0}(J)\ge \nu_{x_0}(\operatorname{in}(J)).
\]
Consequently,
\[
\delta_{\mathbf F}(1)
\ge
\min\{D:x_0^D\in \operatorname{in}(J_{\mathbf F})\}.
\]
\end{corollary}

The affine chart has proof-theoretic meaning, so the homogenizing coordinate should be preserved. Let
\[
P=\{g\in \operatorname{GL}(\operatorname{span}_k(x_0,\ldots,x_n)):g(kx_0)=kx_0\}.
\]
If \(g\in P\), then \(g(x_0)=\lambda x_0\) for some \(\lambda\in k^\times\), and therefore
\[
\nu_{x_0}(gJ)=\nu_{x_0}(J).
\]
Combining this with the preceding corollary gives the family of lower bounds
\[
\delta_{\mathbf F}(1)\ge \nu_{x_0}(\operatorname{in}(gJ_{\mathbf F})),\qquad g\in P.
\]

\section{Finite Domains and Filtered Membership}

Let \(X\subset k^n\) be finite, with vanishing ideal \(I_X\subset R\), coordinate algebra
\[
A_X=R/I_X,
\]
and degree filtration
\[
A_{X,\le D}=\im(R_{\le D}\to A_X).
\]

\begin{definition}[Degree-compatible domain presentation]
A finite presentation \(\Gamma=(\gamma_1,\ldots,\gamma_s)\) of \(I_X\) is degree-compatible if, for every \(D\ge0\) and every \(p\in I_X\cap R_{\le D}\), there are \(u_\ell\in R\) such that
\[
p=\sum_{\ell=1}^s u_\ell\gamma_\ell,\qquad
\deg(u_\ell\gamma_\ell)\le D
\]
for every \(\ell\).
\end{definition}

Finite grids have degree-compatible presentations. If
\[
X=A_1\times\cdots\times A_n
\]
and
\[
u_i(x_i)=\prod_{a\in A_i}(x_i-a),
\]
then \(I_X=(u_1,\ldots,u_n)\), and the presentation is degree-compatible because these univariate polynomials in disjoint variables form a Groebner basis and division does not increase total degree.

\begin{theorem}[Finite-domain filtered membership]
Let \(X\subset k^n\) be finite with degree-compatible presentation \(\Gamma\). Let \(f_1,\ldots,f_m\in R\), and assume
\[
1\in (f_1,\ldots,f_m)+I_X.
\]
For the combined presentation
\[
\mathbf P=(f_1,\ldots,f_m,\Gamma),
\]
one has
\[
\delta_{\mathbf P}(1)
=
\min
\left\{
D:
1\in
\sum_{j=1}^m
\bar f_j\,A_{X,\le D-\deg f_j}
\right\},
\]
where \(\bar f_j\) is the image of \(f_j\) in \(A_X\), and \(A_{X,\le r}=0\) for \(r<0\).
\end{theorem}

\begin{proof}
If
\[
1=\sum_j q_j f_j+\sum_\ell u_\ell\gamma_\ell
\]
is a degree-\(D\) static refutation, then \(\deg q_j\le D-\deg f_j\) whenever \(q_j\ne0\). Passing to \(A_X\) gives
\[
1=\sum_j \bar q_j\bar f_j
\]
with \(\bar q_j\in A_{X,\le D-\deg f_j}\).

Conversely, suppose
\[
1=\sum_j \bar q_j\bar f_j
\]
in \(A_X\), with \(\bar q_j\in A_{X,\le D-\deg f_j}\). Choose representatives \(q_j\in R_{\le D-\deg f_j}\). Then
\[
p=1-\sum_j q_j f_j
\]
lies in \(I_X\) and has degree at most \(D\). Degree-compatibility of \(\Gamma\) writes \(p=\sum_\ell u_\ell\gamma_\ell\) with all products \(u_\ell\gamma_\ell\) of degree at most \(D\). This gives a degree-\(D\) static refutation.
\end{proof}

\section{Single-Equation Reciprocal Formula}

Assume \(X\) has a degree-compatible presentation \(\Gamma\). Let \(f\in R\) have no zeros on \(X\). Then \(\bar f\in A_X\) is invertible. Define
\[
\deg_X(\bar a)=\min\{D:\bar a\in A_{X,\le D}\}.
\]

\begin{theorem}[Reciprocal degree]
For the presentation \(\mathbf P=(f,\Gamma)\),
\[
\delta_{\mathbf P}(1)=\deg f+\deg_X(\bar f^{-1}).
\]
\end{theorem}

\begin{proof}
The finite-domain theorem gives
\[
\delta_{\mathbf P}(1)
=
\min\{D:1\in \bar f A_{X,\le D-\deg f}\}.
\]
Since \(\bar f\) is invertible, this is the least \(D\) such that \(\bar f^{-1}\in A_{X,\le D-\deg f}\).
\end{proof}

This is the central filtered-membership form: for one equation on a finite domain, the multiplier is forced to be the reciprocal function in the finite coordinate algebra, and the exact static degree is its filtered degree plus \(\deg f\).

\section{Boolean Reciprocal Degree}

Let
\[
B_n=(x_i^2-x_i:1\le i\le n),\qquad A_n=R/B_n.
\]
Every class in \(A_n\) has a unique multilinear representative
\[
\sum_{S\subseteq[n]}c_Sx_S,\qquad x_S=\prod_{i\in S}x_i.
\]
Its Boolean degree is
\[
\deg_{\mathrm{bool}}(a)=\max\{|S|:c_S\ne0\}.
\]

\begin{corollary}[Boolean reciprocal formula]
If \(f\in R\) has no zeros on \(\{0,1\}^n\), then
\[
\delta(f,B_n)=\deg f+\deg_{\mathrm{bool}}(f^{-1}),
\]
where \(f^{-1}\) is the inverse of \(f\) in \(A_n\).
\end{corollary}

For a Boolean function \(\phi:\{0,1\}^n\to k\), the coefficient of \(x_1\cdots x_n\) in its multilinear representative is
\[
\sum_{S\subseteq[n]}(-1)^{n-|S|}\phi(\mathbf 1_S).
\]
Thus the top coefficient of \(f^{-1}\) is
\[
C_{\mathrm{top}}(f)
=
\sum_{S\subseteq[n]}
(-1)^{n-|S|}
\frac{1}{f(\mathbf 1_S)}.
\]

\begin{theorem}[Generic maximum on the Boolean cube]
Assume \(k\) is infinite and \(\charac k=0\) or \(\charac k>n\). Let \(d\ge1\). On a nonempty Zariski-open subset of the space of degree-at-most-\(d\) polynomials with no Boolean zeros,
\[
\deg_{\mathrm{bool}}(f^{-1})=n.
\]
Consequently, for generic polynomials of exact degree \(d\) with no Boolean zeros,
\[
\delta(f,B_n)=n+d.
\]
\end{theorem}

\begin{proof}
The condition \(C_{\mathrm{top}}(f)\ne0\) is Zariski-open on the locus where all Boolean values of \(f\) are nonzero. It is enough to show that \(C_{\mathrm{top}}\) is not identically zero. Since \(d\ge1\), consider
\[
f_a=x_1+\cdots+x_n-a,\qquad a\notin\{0,\ldots,n\}.
\]
Then
\[
C_{\mathrm{top}}(f_a)
=
\sum_{j=0}^n(-1)^{n-j}\binom nj\frac1{j-a},
\]
the \(n\)-th finite difference of \(z\mapsto 1/(z-a)\). It equals \(\pm n!/\prod_{j=0}^n(j-a)\), which is nonzero under the characteristic assumptions. The reciprocal formula gives the degree statement.
\end{proof}

\section{Cardinality and Knapsack Contradictions}

Let
\[
t=x_1+\cdots+x_n
\]
and assume \(a\in k\setminus\{0,\ldots,n\}\). Consider
\[
x_i^2-x_i=0\quad(1\le i\le n),\qquad t-a=0.
\]

\begin{theorem}[Cardinality contradiction]
Assume \(\charac k=0\) or \(\charac k>n\). The static refutation degree of this presentation is exactly
\[
n+1.
\]
\end{theorem}

\begin{proof}
The reciprocal formula gives
\[
\delta=1+\deg_{\mathrm{bool}}\bigl((t-a)^{-1}\bigr).
\]
The inverse depends only on the Hamming weight and is the function \(w\mapsto 1/(w-a)\) on \(w=0,\ldots,n\). If it were represented by a polynomial \(q(t)\) of degree \(<n\), then
\[
q(t)(t-a)-1
\]
would be a polynomial of degree at most \(n\) vanishing at the \(n+1\) distinct points \(0,\ldots,n\), hence would be zero. That is impossible because \(t-a\) does not divide \(1\). Interpolation gives a representative of degree at most \(n\), so the degree is exactly \(n\).
\end{proof}

\begin{proposition}[Full multilinear support]
Under the same assumptions, the multilinear inverse \(q=(t-a)^{-1}\in A_n\) has every squarefree monomial with nonzero coefficient. Hence
\[
|\Supp(q)|=2^n.
\]
Moreover every polynomial representative of \(q\) has at least \(2^n\) standard monomials before reduction modulo \(B_n\).
\end{proposition}

\begin{proof}
By symmetry, the coefficient of a squarefree monomial of degree \(m\) is
\[
c_m=\sum_{j=0}^m(-1)^{m-j}\binom mj\frac1{j-a}
=\pm\frac{m!}{\prod_{j=0}^m(j-a)}.
\]
This is nonzero for \(0\le m\le n\). Reduction modulo \(B_n\) maps each standard monomial to one squarefree monomial, so any representative must contain at least the support appearing in the multilinear normal form.
\end{proof}

This support statement is a normal-form size statement. It is compatible with the fact that the same multiplier is a univariate polynomial in \(t\) of degree \(n\).

The same argument applies to weighted knapsack equations. If
\[
L(x)=a+\sum_{i=1}^n c_i x_i
\]
has no Boolean zeros, then
\[
\delta(L,B_n)=1+\deg_{\mathrm{bool}}(L^{-1}),
\]
and for a nonempty Zariski-open set of coefficients this degree is \(n+1\).

\section{Equivariant Macaulay Maps}

Let \(G\) be a linearly reductive group over \(k\). Assume \(G\) acts linearly on \(\operatorname{span}_k(x_1,\ldots,x_n)\) and fixes \(x_0\).

A homogeneous \(G\)-equivariant presentation consists of finite-dimensional \(G\)-modules \(E_a\), degrees \(d_a\), and \(G\)-equivariant maps
\[
\alpha_a:E_a\to S_{d_a}.
\]
Let
\[
J_\alpha=\left(\sum_a \alpha_a(E_a)\right)\subset S.
\]
For \(D\ge0\), define
\[
M_D(\alpha)=\bigoplus_a S_{D-d_a}\otimes E_a,
\]
with \(S_j=0\) for \(j<0\), and the Macaulay map
\[
\mu_D:M_D(\alpha)\to S_D,\qquad m\otimes e\mapsto m\alpha_a(e).
\]
Then \((J_\alpha)_D=\im\mu_D\), and \(\mu_D\) is \(G\)-equivariant.

\begin{theorem}[Invariant Macaulay reduction]
Let \(H\in S_D^G\). Then
\[
H\in (J_\alpha)_D
\quad\Longleftrightarrow\quad
H\in \im\bigl(\mu_D^G:M_D(\alpha)^G\to S_D^G\bigr).
\]
In particular,
\[
x_0^D\in J_\alpha
\quad\Longleftrightarrow\quad
x_0^D\in \im(\mu_D^G).
\]
\end{theorem}

\begin{proof}
One implication is immediate. Conversely, if \(H=\mu_D(u)\), apply the Reynolds projection \(\rho:M_D(\alpha)\to M_D(\alpha)^G\). Functoriality gives
\[
\mu_D(\rho(u))=\rho(\mu_D(u))=\rho(H)=H.
\]
\end{proof}

Thus invariant static proof search is compressed exactly to invariant Macaulay maps.

\section{Isotypic Obstructions}

Let
\[
S_D=\bigoplus_\lambda S_D[\lambda],
\qquad
M_D(\alpha)=\bigoplus_\lambda M_D(\alpha)[\lambda]
\]
be the isotypic decompositions. Since \(\mu_D\) is \(G\)-equivariant, it preserves isotypic type.

\begin{theorem}[Isotypic membership criterion]
If \(H\in S_D[\lambda]\), then
\[
H\in (J_\alpha)_D
\quad\Longleftrightarrow\quad
H\in \mu_D(M_D(\alpha)[\lambda]).
\]
In particular, if \(M_D(\alpha)[\lambda]=0\) and \(H\ne0\), then \(H\notin (J_\alpha)_D\).
\end{theorem}

\begin{proof}
If \(H=\mu_D(u)\), decompose \(u=\sum_\eta u_\eta\) by isotypic type. Then
\[
H=\sum_\eta \mu_D(u_\eta),
\]
with \(\mu_D(u_\eta)\in S_D[\eta]\). Since \(H\) has type \(\lambda\), all other components vanish, and \(H=\mu_D(u_\lambda)\).
\end{proof}

First occurrence of representations in the Macaulay source therefore gives degree lower bounds. If all generators are \(G\)-invariant of degrees \(d_i\), a type-\(\lambda\) consequence cannot be derived in degree \(D\) unless \(\lambda\) occurs in some \(S_{\le D-d_i}\).

\section{Symmetric Boolean Systems}

Assume \(\charac k=0\) or \(\charac k>n\). Let \(S_n\) act on
\[
A_n=k[x_1,\ldots,x_n]/(x_i^2-x_i)
\]
by permuting variables, and put \(t=x_1+\cdots+x_n\).

\begin{lemma}[Invariant Boolean quotient]
There is a filtered algebra isomorphism
\[
A_n^{S_n}
\cong
k[t]/(\chi_n(t)),
\qquad
\chi_n(t)=\prod_{a=0}^n(t-a),
\]
where the degree filtration corresponds to ordinary degree in \(t\), truncated at \(n\).
\end{lemma}

\begin{proof}
A symmetric Boolean function depends only on Hamming weight \(t\in\{0,\ldots,n\}\), so the invariant algebra is the coordinate algebra of this finite set. The filtration statement follows because the elementary symmetric function \(e_j\) restricts to \(\binom tj\), and \(j!\) is invertible for \(0\le j\le n\).
\end{proof}

\begin{theorem}[Exact univariate reduction]
Let \(f_1,\ldots,f_m\in k[x_1,\ldots,x_n]^{S_n}\), and consider the Boolean presentation
\[
\mathbf P=(f_1,\ldots,f_m,\ x_i^2-x_i\ (1\le i\le n)).
\]
Let \(\bar f_j(t)\) be the image of \(f_j\) in \(k[t]/(\chi_n(t))\). Then \(\delta_{\mathbf P}(1)\) is the least \(D\) such that
\[
1\in
\sum_{j=1}^m
\bar f_j(t)\,k[t]_{\le D-\deg f_j}
\quad\text{inside }k[t]/(\chi_n(t)).
\]
\end{theorem}

\begin{proof}
The finite-domain theorem gives filtered membership in \(A_n\). Since the target \(1\) and the data are invariant, Reynolds reduction allows restriction to \(A_n^{S_n}\). The invariant Boolean quotient identifies this with the displayed univariate problem.
\end{proof}

\section{Block Symmetry and Finite Grids}

Let
\[
[n]=B_1\sqcup\cdots\sqcup B_r,\qquad |B_\ell|=n_\ell,
\]
and let \(G=S_{n_1}\times\cdots\times S_{n_r}\) act by permuting variables inside each block. Define
\[
t_\ell=\sum_{i\in B_\ell}x_i.
\]
If \(\charac k=0\) or \(\charac k>\max_\ell n_\ell\), then
\[
A_n^G
\cong
k[t_1,\ldots,t_r]/
(\chi_{n_1}(t_1),\ldots,\chi_{n_r}(t_r)),
\]
with total-degree filtration in the orbit parameters.

Consequently, a block-symmetric Boolean system
\[
F_j(t_1,\ldots,t_r)=0
\]
has static refutation degree equal to the least \(D\) such that
\[
1\in
\sum_j \bar F_j\,k[t_1,\ldots,t_r]_{\le D-\deg F_j}
\]
inside the finite-grid coordinate algebra above. For one equation \(F=0\) with no zeros on the grid,
\[
\delta=\deg F+\deg_{\mathrm{grid}}(F^{-1}).
\]

Let \(N=n_1+\cdots+n_r\). On coefficient spaces where a top grid coefficient of \(F^{-1}\) is not identically zero, generic degree-\(d\) block-symmetric contradictions have degree
\[
N+d.
\]
For affine-linear \(F\), nonvanishing is witnessed by \(F=t_1+\cdots+t_r-a\), with \(a\notin\{0,\ldots,N\}\), using iterated finite differences of \(1/(z-a)\).

\section{Dual Linear Functionals}

Let
\[
V_D=\sum_i f_i R_{\le D-\deg f_i}\subset R_{\le D}.
\]
Then \(\delta_{\mathbf F}(1)>D\) if and only if \(1\notin V_D\), equivalently there exists a linear functional
\[
\Lambda:R_{\le D}\to k
\]
such that
\[
\Lambda(1)=1,\qquad
\Lambda(f_iq)=0
\]
for every \(i\) and every \(q\in R_{\le D-\deg f_i}\).

Over a finite domain, this dual functional may be taken on the filtered coordinate algebra.

\begin{proposition}[Reciprocal dual form]
Let \(X\) be a finite grid and let \(f\) be nonzero on \(X\). If
\[
\deg_X(f^{-1})>r,
\]
then there exists a linear functional
\[
\Lambda:A_{X,\le r+\deg f}\to k
\]
such that
\[
\Lambda(1)\ne0,\qquad
\Lambda(\bar f\,A_{X,\le r})=0.
\]
\end{proposition}

\begin{proof}
Since \(f^{-1}\notin A_{X,\le r}\), the subspace \(\bar f A_{X,\le r}\) does not contain \(1\). Finite-dimensional linear algebra separates \(1\) from that subspace.
\end{proof}

If a linearly reductive group acts and all data are invariant, such functionals may be averaged to invariant functionals.

\section{First Occurrence for Symmetric Groups}

Assume \(\charac k=0\). Let \(S_n\) act on \(R=k[x_1,\ldots,x_n]\). For a partition \(\lambda\vdash n\), let \(S^\lambda\) be the Specht module. The first total degree in which \(S^\lambda\) occurs in \(R\) is
\[
b(\lambda)=\sum_i(i-1)\lambda_i,
\]
the lowest degree appearing in the fake-degree polynomial.

\begin{proposition}[First-occurrence bound]
Let \(f_1,\ldots,f_m\in R^{S_n}\) be invariant generators of degrees \(d_1,\ldots,d_m\). If a homogeneous component of a consequence has type \(\lambda\), then no static derivation of that component in degree \(D\) can occur unless
\[
D\ge \min_i(d_i+b(\lambda)).
\]
\end{proposition}

\begin{proof}
Since the \(f_i\) are invariant, \(q_if_i\) has the same \(S_n\)-type as \(q_i\). In degree \(D\), the multiplier \(q_i\) has degree at most \(D-d_i\). If \(D-d_i<b(\lambda)\) for every \(i\), then the Macaulay source has no \(\lambda\)-isotypic component, and the isotypic membership criterion applies.
\end{proof}

For the sign representation \(\lambda=(1^n)\),
\[
b(\lambda)=\binom n2.
\]
Thus invariant generators of minimum degree \(d_{\min}\) cannot produce an alternating component before degree
\[
d_{\min}+\binom n2.
\]

\section{Position Among Algebraic Proof Measures}

The results above concern static algebraic derivation degree:
\[
h=\sum_i g_i f_i,\qquad \max_i\deg(g_if_i).
\]
They apply directly to Hilbert--Nullstellensatz degree and bounded static ideal membership. They give exact filtered-degree and equivariant reductions for such identities.

They also explain a precise bridge to functional lower-bound methods. For a single equation \(f=0\) over a finite domain, the multiplier is the reciprocal function \(1/f\) in the coordinate algebra. Lower bounds for static degree are exactly lower bounds on the filtered degree of this reciprocal. Stronger proof systems, such as IPS fragments or dynamic Polynomial Calculus size measures, require additional circuit-complexity or derivational arguments, but the finite-domain reciprocal identity is the static-degree core of that phenomenon.

\section{Summary}

The main conclusions are:
\begin{enumerate}[label=(\roman*)]
\item for a chosen presentation,
\[
\delta_{\mathbf F}(h)=\min\{D:x_0^{D-\deg h}h^h\in J_{\mathbf F}\};
\]
\item the excess degree is the \(x_0\)-torsion order in \((J_{\mathbf F}:x_0^\infty)/J_{\mathbf F}\);
\item over degree-compatible finite domains,
\[
\delta(1)=
\min\left\{D:
1\in\sum_j \bar f_j A_{X,\le D-\deg f_j}
\right\};
\]
\item for one equation with no finite-domain zeros,
\[
\delta=\deg f+\deg_X(f^{-1});
\]
\item generic degree-\(d\) Boolean contradictions have degree \(n+d\);
\item the cardinality contradiction has exact degree \(n+1\) and a full-support multilinear reciprocal multiplier;
\item invariant static proof search under a linearly reductive group is exactly compressed by invariant Macaulay maps;
\item symmetric Boolean systems reduce exactly to univariate filtered membership in
\[
k[t]/\prod_{a=0}^n(t-a);
\]
\item non-invariant consequences admit first-occurrence lower bounds from isotypic representation theory.
\end{enumerate}

In the language of access systems, \(\delta_{\mathbf F}(h)\) is the certification cost associated with the chosen polynomial presentation. The homogenized ideal \(J_{\mathbf F}\) records which consequences are visible at each degree level, while the \(x_0\)-torsion quotient records the delay before a true affine consequence becomes visible in the filtered presentation. Thus static proof degree is filtered membership: a bounded certificate exists exactly when the homogenized consequence appears in the corresponding finite degree slice.

\begin{thebibliography}{99}

\bibitem{BureshClegg}
J. Buresh-Oppenheim, M. Clegg, R. Edmonds, R. Impagliazzo, and T. Pitassi.
\newblock Homogenization and the Polynomial Calculus.
\newblock \emph{Computational Complexity} 11 (2002), 91--108.

\bibitem{CleggEdmondsImpagliazzo}
M. Clegg, J. Edmonds, and R. Impagliazzo.
\newblock Using the Groebner Basis Algorithm to Find Proofs of Unsatisfiability.
\newblock In \emph{Proceedings of STOC 1996}, 174--183.

\bibitem{ForbesEtAl}
M. A. Forbes, A. Shpilka, I. Tzameret, and A. Wigderson.
\newblock Proof Complexity Lower Bounds from Algebraic Circuit Complexity.
\newblock \emph{Theory of Computing} 17 (2021).

\bibitem{GrochowPitassi}
J. A. Grochow and T. Pitassi.
\newblock Circuit Complexity, Proof Complexity, and Polynomial Identity Testing: The Ideal Proof System.
\newblock \emph{Journal of the ACM} 65(6) (2018).

\bibitem{HakoniemiLimayeTzameret}
T. Hakoniemi, N. Limaye, and I. Tzameret.
\newblock Functional Lower Bounds in Algebraic Proofs: Symmetry, Lifting, and Barriers.
\newblock arXiv:2412.20114, 2024.

\bibitem{SombraHilbert}
M. Sombra.
\newblock Bounds for the Hilbert Function of Polynomial Ideals and for the Degrees in the Nullstellensatz.
\newblock \emph{Journal of Pure and Applied Algebra} 117/118 (1997), 565--599.

\bibitem{SombraSparse}
M. Sombra.
\newblock A Sparse Effective Nullstellensatz.
\newblock \emph{Advances in Applied Mathematics} 22 (1999), 271--295.

\bibitem{MayrMeyer}
E. W. Mayr and A. R. Meyer.
\newblock The Complexity of the Word Problems for Commutative Semigroups and Polynomial Ideals.
\newblock \emph{Advances in Mathematics} 46 (1982), 305--329.

\bibitem{BayerStillman}
D. Bayer and M. Stillman.
\newblock A Criterion for Detecting \(m\)-Regularity.
\newblock \emph{Inventiones Mathematicae} 87 (1987), 1--11.

\bibitem{Eisenbud}
D. Eisenbud.
\newblock \emph{Commutative Algebra with a View Toward Algebraic Geometry}.
\newblock Springer, 1995.

\bibitem{CoxLittleOShea}
D. Cox, J. Little, and D. O'Shea.
\newblock \emph{Ideals, Varieties, and Algorithms}.
\newblock Springer.

\bibitem{FultonHarris}
W. Fulton and J. Harris.
\newblock \emph{Representation Theory: A First Course}.
\newblock Springer, 1991.

\bibitem{Sagan}
B. E. Sagan.
\newblock \emph{The Symmetric Group: Representations, Combinatorial Algorithms, and Symmetric Functions}.
\newblock Springer.

\bibitem{Macdonald}
I. G. Macdonald.
\newblock \emph{Symmetric Functions and Hall Polynomials}.
\newblock Oxford University Press.

\bibitem{SamSnowden}
S. V. Sam and A. Snowden.
\newblock GL-Equivariant Modules over Polynomial Rings in Infinitely Many Variables.
\newblock \emph{Transactions of the American Mathematical Society} 368 (2016).

\end{thebibliography}

\end{document}
